\section{Game-Theoretic Analysis}
\label{sec:game}

\subsection{Validator Participation Game}

We define the participation game $\Gamma = (\mathcal{N}, \{A_i\}, \{U_i\})$ where:
\begin{itemize}
    \item $\mathcal{N} = \{1, \ldots, N\}$: potential validators
    \item $A_i = \{0\} \cup [s^{\min}_{\text{val}}, s_i^{\max}]$: action space (stake 0 or $\geq s^{\min}_{\text{val}}$)
    \item $U_i$: utility function
\end{itemize}

\begin{definition}[Validator Utility]
\begin{equation}
    U_i(s_i, s_{-i}) = \underbrace{r_i^{\text{self}}(s_i, s_{-i})}_{\text{self-stake reward}} + \underbrace{r_i^{\text{commission}}(s_i, s_{-i})}_{\text{commission income}} - \underbrace{h_i}_{\text{hardware cost}} - \underbrace{s_i \cdot r_{\text{opp}}}_{\text{opportunity cost}} - \underbrace{\mathbb{E}[\text{slash}_i]}_{\text{expected slashing}}
\end{equation}
\end{definition}

\subsection{Honest Validation as Dominant Strategy}

\begin{theorem}[Incentive Compatibility]\label{thm:ic}
For a validator $v_i$ with $s_i < W/3$, honest participation is a strictly dominant strategy when:
\begin{equation}
    r_i^{\text{total}} > h_i + s_i \cdot r_{\text{opp}} + \max_{\text{attack}} \left(\pi_{\text{attack}} - p_{\text{slash}} \cdot s_i \cdot p_{\text{detect}}\right)
\end{equation}
where $\pi_{\text{attack}}$ is the maximum extractable value from any attack and $p_{\text{detect}}$ is the detection probability.
\end{theorem}

\begin{proof}
The utility from honest behavior is $U^H_i = r_i^{\text{total}} - h_i - s_i \cdot r_{\text{opp}}$. The utility from any Byzantine strategy is:
\begin{equation}
    U^B_i = r_i^{\text{total}} + \pi_{\text{attack}} \cdot (1 - p_{\text{detect}}) - p_{\text{slash}} \cdot s_i \cdot p_{\text{detect}} - h_i - s_i \cdot r_{\text{opp}}
\end{equation}

For equivocation, $p_{\text{detect}} = 1$ (cryptographically provable), so:
\begin{equation}
    U^B_i - U^H_i = \pi_{\text{attack}} \cdot 0 - 0.05 \cdot s_i < 0
\end{equation}

For coordinated attacks requiring $\geq W/3$ stake, $v_i$ alone (with $s_i < W/3$) cannot mount the attack. Collusion with others introduces the free-rider problem: each participant bears the full slashing risk but shares the attack profit, making defection (honest behavior) individually rational.
\end{proof}

\subsection{Delegation Market Equilibrium}

\begin{theorem}[Delegation Equilibrium]
In a market with $n$ validators and $M$ total delegatable tokens, the equilibrium delegation allocation satisfies:
\begin{equation}
    \frac{s_d^*}{M} = \frac{(1 - c_i) \cdot u_i}{\sum_j (1 - c_j) \cdot u_j} \quad \text{for delegators of } v_i
\end{equation}
That is, delegation flows proportionally to the after-commission, uptime-adjusted return.
\end{theorem}

\begin{proof}
A rational delegator maximizes expected return $(1 - c_i) \cdot u_i \cdot R / W$. In equilibrium, all validators offer the same marginal return to delegators. If validator $i$ offers higher marginal return, it attracts delegation until marginal returns equalize (diminishing returns from increased $W$).
\end{proof}

\subsection{Entry and Exit Dynamics}

\begin{proposition}[Free Entry Equilibrium]
The equilibrium number of validators $n^*$ satisfies:
\begin{equation}
    n^* = \left\lfloor \frac{R_t \cdot c^{\min}}{h^{\text{avg}} + s^{\min}_{\text{val}} \cdot r_{\text{opp}}} \right\rfloor
\end{equation}
where $h^{\text{avg}}$ is the average hardware cost per epoch.
\end{proposition}

\begin{proof}
Entry occurs when $U_i > 0$; exit occurs when $U_i < 0$. At the margin, the least profitable validator earns zero profit: commission income equals costs. With $n$ validators each earning approximately $R_t / n$ in base rewards, the marginal validator earns $c^{\min} \cdot R_t / n$ in commission. Setting this equal to $h + s^{\min}_{\text{val}} \cdot r_{\text{opp}}$ and solving for $n$ yields the result.
\end{proof}

